\hmsubsection{Calibration Subsystem}{UMA}{}\label{sec:calibration}

The geometric model used by the inverse kinematics solver depends on 
several parameters that cannot be measured analytically: the 
sag coefficients that describe the arm's elastic deformation, the 
pixel-to-centimetre transformation that maps camera images to 
workspace coordinates, and the joint configuration at which the 
camera observes the workspace. These parameters are determined 
through a set of dedicated calibration procedures that are run 
once during initial system setup and re-run whenever the physical 
configuration of the arm changes.

Calibration is treated in the present project as a software 
subsystem rather than a one-time setup activity. Each procedure 
is implemented as a standalone Python script that produces a JSON 
file containing the calibrated parameters, and each parameter file 
is loaded automatically by the corresponding runtime component. 
This separation between calibration and runtime keeps the source 
code free of hand-tuned constants and allows the system to be 
recalibrated without recompilation or modification of the source. 
Three procedures directly support the inverse kinematics solver: 
sag calibration (Section~\ref{sec:sag-calibration}), touch calibration 
(Section~\ref{sec:touch-cal}), and scan-pose calibration (Section~\ref{sec:scan-pose-calibration}).

\hmsubsubsection{Sag Calibration}{UMA}{}\label{sec:sag-calibration}

The sag calibration procedure determines the coefficients of the 
empirical droop model described in Section~\ref{sec:sag}. The procedure 
disables the runtime sag compensation in the solver and commands 
the arm to a fixed height above the workspace surface at six 
horizontal reach distances between 6 and 36~cm. The calibration 
run that produced the present coefficients used a commanded 
height of 2.0~cm; the script's default has since been raised to 
10.0~cm to avoid desk contact at long reaches. At each reach 
distance, the operator measures the actual height of the gripper 
tip with a ruler and enters the value into the script. The 
script computes the difference between the commanded and the 
measured height at each point, fits both a linear and a 
quadratic regression model to the resulting droop measurements, 
and writes the fitted coefficients to the file 
\texttt{sag\_calibration.json}. The quadratic model is selected 
automatically when its root mean square error is at least 20\% 
lower than that of the linear model; otherwise the operator may 
override the selection manually.

The fitted coefficients are loaded automatically by the inverse 
kinematics solver at startup. If the file is missing or contains 
invalid data, the solver falls back to a set of default 
coefficients embedded in the source code, which produce 
approximately correct compensation but are inferior to a fresh 
calibration of the actual arm. In the deployed configuration, 
the sag model is not exercised during normal operation because 
all pick operations and route waypoints use per-point surface 
heights from the touch calibration procedure 
(Section~\ref{sec:touch-cal}), which bypass the polynomial correction via the 
solver's \texttt{skip\_sag} flag. The sag model therefore serves 
as a fallback for deployments without touch calibration data. 
The complete procedure takes approximately 10--20 minutes and 
must be repeated whenever the arm is mechanically modified, such 
as after replacing a motor or adjusting the bracket geometry.

\hmsubsubsection{Touch Calibration}{UMA}{}\label{sec:touch-cal}

The touch calibration procedure determines the perspective 
transformation that maps pixel coordinates in the camera image 
to physical centimetres in the workspace coordinate frame. The 
procedure was designed as a replacement for an earlier 
ruler-based calibration method that introduced significant 
measurement error and required the operator to mentally align 
two distinct coordinate frames.

\textbf{Method.} The operator places between four and twenty 
balls on the workspace surface, with the precise number being a 
trade-off between calibration time and the statistical 
robustness of the resulting transformation. The arm moves to 
the scan pose and the camera detects the ball positions in 
pixel coordinates, averaging detections across forty frames to 
suppress single-frame noise. The operator then drives the arm 
manually so that the gripper physically touches each ball in 
turn. When the gripper is centred on a ball, the operator 
confirms the position and the script records the corresponding 
workspace coordinates directly from the inverse kinematics 
solver, eliminating any ruler measurement.

When more than four balls are used, the perspective 
transformation is computed by the random sample consensus 
algorithm \cite{fischler1981} rather than by direct 
four-point construction. The random sample consensus algorithm 
treats individual measurements as inliers or outliers and fits 
a transformation that minimises the residual error across the 
inlier set. This produces a more robust transformation than 
the direct construction, which gives equal weight to every 
measurement regardless of its quality.

\textbf{Limp-mode operation.} An additional mode is provided 
in which the operator drives the arm by hand rather than 
through the keyboard. In this mode, the torque on motors~1 
through~4 is disabled, allowing the arm to be physically 
guided to each ball. When the gripper reaches the desired 
position, the firmware reports the current encoder values back 
to the host, and the forward kinematics described in 
Section~\ref{sec:fk-validation} converts these values into a workspace position. 
The torque is then immediately re-enabled with the goal 
positions set to the recorded values, preventing the arm from 
falling under its own weight.

The limp-mode procedure is more accurate than keyboard-driven 
calibration for points near the workspace edges, since the 
operator can position the gripper directly without iterating 
between coordinate axes. A subsequent fine-tuning phase, in 
which the arm is driven by inverse kinematics to a clearance 
height above each captured point and the operator adjusts the 
final position with the keyboard, refines the calibration data 
without requiring the operator to support the arm's weight.

\textbf{Per-point auxiliary data.} In addition to the pixel 
and workspace coordinates used for the perspective 
transformation, the procedure records two auxiliary values at 
each calibration point: the height at which the gripper 
contacts the surface, and the wrist-tilt offset that produces 
a level gripper at that position. These values are used at 
runtime to interpolate corrections specific to the position 
of the target. The recorded surface height supersedes the 
sag model during gripping operations, since direct 
measurement at known positions is more accurate than a 
reach-based regression. To prevent the two compensations from 
being applied simultaneously, the inverse kinematics solver 
accepts a flag that disables the sag model when surface-height 
data is available for the requested position.

\textbf{Output.} The procedure writes the perspective 
transformation, the per-point pixel and workspace 
coordinates, the recorded surface heights, and the recorded 
wrist offsets to the file 
\texttt{homography\_calibration.json}. This file is loaded by 
the vision bridge component, which uses the perspective 
transformation to convert ball detections into workspace 
positions, and by the configuration module, which provides 
the per-point interpolation functions to the inverse 
kinematics solver. The complete procedure takes approximately 
five to ten minutes.

\hmsubsubsection{Scan Pose Calibration}{UMA}{}\label{sec:scan-pose-calibration}

The scan pose is the joint configuration at which the arm 
positions the camera over the workspace for ball detection. 
Because the camera is mounted on the wrist of the arm, its 
orientation and position depend on the joint angles of the 
arm rather than being fixed in the workspace. The perspective 
transformation produced by the touch calibration procedure 
is therefore valid only at one specific joint configuration, 
which must be determined and recorded so that the arm can 
return to it reliably before each ball detection.

A fixed scan pose was selected for this purpose, in 
preference to a continuously updated transformation that 
would track the camera's position as the arm moves. The fixed 
approach was chosen for two reasons. First, the alternative 
requires accurate forward kinematics for the wrist-mounted 
camera frame, which is difficult to achieve in the presence 
of sag and would need to be recalibrated alongside the arm 
itself. Second, the gripper enters the camera's field of 
view during the final approach to a ball, so visual feedback 
during the approach is unreliable regardless of the chosen 
method.

The scan pose calibration procedure is implemented as an 
interactive script that opens the camera feed and allows the 
operator to adjust each joint with keyboard controls until 
the camera looks directly down at the workspace and the 
entire reachable area is visible in the frame without the 
gripper obscuring the view. The selected joint angles are 
written to the configuration file \texttt{src/config/arm.py}, 
where they are accessible to the runtime as a constant. The 
procedure takes approximately five to ten minutes and must 
be repeated whenever the camera is remounted or the arm's 
mechanical configuration changes.

\textbf{Pose stability requirement.} The reliability of the 
scan pose depends on the arm's ability to return to it 
repeatably. The Dynamixel servos used in the project provide 
sub-degree positional repeatability under nominal load, but the scan
pose must also avoid imposing unnecessary static load on the elbow
motor. As described in Section~\ref{sec:thermal}, a reduced elbow
current limit was evaluated but removed from the deployed system because
it did not measurably reduce the thermal load. The deployed scan pose
therefore uses mechanical support instead: the L\textsubscript{2}
structure rests on the L\textsubscript{1} piece while preserving the
camera's view of the workspace. This reduces the holding torque required
from M3 during idle scanning while maintaining a repeatable camera
configuration. The scan pose therefore relies on mechanical stability,
repeatable joint positioning, and verification of the camera calibration
before test sessions rather than on a reduced holding-current setting.
This reduces the idle holding load at \texttt{SCAN\_POSE}, but it does
not remove thermal loading caused by movements between poses. Any camera
displacement caused by joint droop or remounting invalidates the
perspective transformation and requires the scan pose calibration to be
repeated.
